\chapter{$G$-Sets and $G$-Graphs}

\section{Permutation Groups}

% Definition of permutation group:
\begin{definition}[Permutation group]
A \emph{permutation group} is a triple $(G, \Omega, \rho)$
where $G$ is a group, $\Omega$ is a set and $\rho$
is a homomorphism:
$$ \rho : G \to \Aut(\Omega). $$
We say that $G$ {\it acts on} $\Omega$ as a group of permutations.
The action is said to be \emph{faithful} if $\ker(\rho) = \{1\}$.
\end{definition}

% Remark on notation:

We will usually neglect to mention the homomorphism $\rho$ explicitely,
and speak of the permutation group $(G, \Omega)$.
We write $\alpha^{\rho(g)}$ or just $\alpha^g$ to indicate
the action of a permutation $g \in G$ on a point $\alpha \in \Omega$. 
This allows us to compose permutations as 
$\alpha^{(gh)} = (\alpha^{g})^{h}$
rather than $(gh)(\alpha) = h(g(\alpha))$.

If $\Delta$ is a subset of $\Omega$ and $g$ is an element of $G$, then we write
$\Delta^g = \{ \alpha^g : \alpha \in \Delta \}$
to indicate the image of $\Delta$ under the action of $g$.

\hfill \break
Sometimes we will call $(G,\Omega)$ a \emph{representation} 
of $G$ as a group of permutations of $\Omega$, rather than a permutation group.
If $|\Omega| = n$ then we say that $(G,\Omega)$ is a permutation
representation of $G$ of degree n.

% Definition: permutation equivalence:
\begin{definition}[Permutation Equivalence]
Two permutation groups $(G_1, \Omega_1)$ and
$(G_2, \Omega_2)$ are said to be \emph{permutation equivalent}
if there is an isomorphism $\varphi : G_1 \to G_2$
and a bijection $\eta : \Omega_1 \to \Omega_2$ 
such that for every $g \in G_1$ the following diagram commutes:

% Commutative diagram using xypic:
\centerline{
 \xymatrix{
 \ar[d]_{\eta} \ar[r]^g \Omega_1 & \Omega_1 \ar[d]_{\eta} \\
 \ar[r]_{\varphi(g)} \Omega_2 & \Omega_2
 }
}
\end{definition} 

Permutation equivalence is an equivalence relation
on the set of permutation representations of given group.
If we relax the condition that $\eta$ is bijective 
we obtain a {\it permutation homomorphism}.
%For a fixed group $G$,
%the set of permutation representations of $G$, 
%together with the permutation homomorphisms form a category,
%which we denote by $G$-Set.

% Definition of transitivity
\begin{definition}[Transitive Permutation Group]
A permutation group  $(G, \Omega)$ is said to be \emph{transitive}
if for any $\alpha, \beta \in \Omega$ there exist a $ g \in G $
such that $\alpha^g = \beta$
\end{definition}

% Definition of orbit
\begin{definition}[Orbit]
If $(G, \Omega)$ is a permutation group and $\alpha \in \Omega$ 
then the \emph{orbit of $\alpha$ under the action of $G$} is:
$$ \alpha^G = \{ \alpha^x : x \in G \} $$
\end{definition}

For a transitive permutation group $\alpha^G = \Omega$.
Otherwise $\Omega$ maybe be partitioned into orbits:
$ \Omega = \coprod_{i = 1..n} \Omega_i $
such that for each $i$, the permutation group $(G,{\Omega_i})$ is transitive.

% category theory
%\hfill \break
%In category theoretic language, the transitive representations
%of a group $G$ are precisely the objects of $G$-Set
%which have no proper subobjects.

\section{Graphs}

The most general definition of a {\it graph} is a pair of sets $(V, A)$
with $A \subseteq V \times V$.
The elements of $V$ are refered to as the {\it vertices} of the graph
and the elements of $A$ are refered to as {\it arcs}.
A {\it simple graph} is a graph such that:
$$ (\alpha,\beta) \in A \Longleftrightarrow (\beta, \alpha) \in A
\mbox{ for all } \alpha, \beta \in V.$$ 

A graph that is not simple is said to be {\it directed}.
The term ``digraph'' is also frequenty used.
An {\it edge} of a graph is an unordered pair $\{\alpha,\beta\}$
where $(\alpha,\beta)$ is an arc and $\alpha \neq \beta$.
A {\it loop} is an arc of the form $(\alpha,\alpha)$.
In this paper we will be mostly interested
in simple graphs without loops
and shall use the term ``graph'' to mean
``simple graph without loops''. We give a formal definition for ease of reference.

\begin{definition} [Graph]
A {\it graph} is a pair of sets $(V,A)$ such that:
$$A \subseteq (V \times V) \diagdown \{ (\alpha,\alpha) : \alpha \in V \}$$

and: 
$$(\alpha,\beta) \in A \Longleftrightarrow (\beta,\alpha) \in A
\mbox{ for all } \alpha, \beta \in V.$$
\end{definition}

When more than one graph is being discussed, we shall write
$\Gamma = (V \Gamma, A \Gamma)$ or $\Sigma = (V \Sigma, A \Sigma)$
so that it is clear to which graph each set of vertices and arcs belongs.
If $\alpha$ is a vertex of some graph $\Gamma$,
then the {\it neighbourhood} of $\alpha$ in $\Gamma$
is the set $\Gamma(\alpha) = \{ \beta \in V \Gamma : (\alpha,\beta) \in A \Gamma \}$.
If $\alpha$ is a vertex of the graph $\Sigma$
then the neighbourhood of $\alpha$ in $\Sigma$ is denoted by $\Sigma(\alpha)$.
We write $E \Gamma$ (or $E \Sigma$) to denote the set of edge of $\Gamma$ (or $\Sigma$)
and $N \Gamma$ (or $N \Sigma$) to denote the set of neighbourhoods.

\begin{definition} [Automorphism of a Graph]
An automorphism of a graph $\Gamma$ is a bijection
$\varphi : V \Gamma \to V \Gamma$ such that:
$$ (\alpha,\beta) \in A \Gamma \Longleftrightarrow
(\varphi(\alpha),\varphi(\beta)) \in A \Gamma \mbox{ for all } \alpha,\beta \in V \Gamma.$$ 
\end{definition}

The autmorphisms of a graph form a group which we denote by $\Aut(\Gamma)$.
If $\rho : G \to \Aut(\Gamma)$ is a homomorphism
then $(G,V \Gamma,\rho)$ is a permutation group
and we say that $G$ acts on $\Gamma$ as a group of automorphisms.

\begin{definition} [Vertex-transitive Graph]
A vertex transitive graph is a triple $(G,\Gamma,\rho)$
where $\rho : G \to \Aut(\Gamma)$ is a homomorphism
and $(G,\Gamma,\rho)$ is transitive.
\end{definition}

As with permutation groups, we will often neglect to mention $\rho$
explicitely. We shall also sometimes say that $\Gamma$ is a $G$-vertex
transitive graph rather than $(G,\Gamma)$ is a vertex transitive graph.
If we say simply that $\Gamma$ is a vertex transitive graph
then we mean that $(\Aut(\Gamma),\Gamma)$ 
is an vertex transitive graph.
%Need to define the valency of a graph.
%Also complete graph and bipartite graph.

\section{Local Properties}

\begin{definition}[Stabilizer]
Let $(G, \Omega)$ be a transitive permutation group,
and let $\alpha$ be any point of $\Omega$.
The \emph{point stabilizer} of $\alpha$ is:
$$ G_a = \{ x \in G : \alpha^x = \alpha  \} $$
\end{definition}

\begin{lemma}
$G_{\alpha^g} = g^{-1} G_{\alpha} g$
for any $g \in G$. 
\begin{proof}
Suppose that $x \in G$ is such that $(\alpha^g)^x = \alpha^g$.
It follows that $\alpha^{gxg^{-1}} = \alpha$
and so $gxg^{-1} \in G_\alpha$. That is $x \in g^{-1}G_\alpha g$.
Conversely suppose that $x \in g^{-1}G_\alpha g$.
Then $gxg^{-1} \in G_\alpha$ and so $\alpha^{gxg^{-1}} = \alpha$.
That is $\alpha^{gx} = \alpha$, so $x \in G_{\alpha^g}$.
\end{proof}
\end{lemma}

Since $(G,\Omega)$ is transitive, every element of $\Omega$
may be expressed in the form $\alpha^g$ for some $g \in G$.
Thus lemma 1 tells us that the point stabilizers of a transitive
representation of $G$ form a family of conjugate subgroups of $G$.

If $G$ acts on $\Gamma$ as a group of automorphisms
and $G_\alpha$ is the stabilizer of the vertex $\alpha$
then $(G_\alpha, \Gamma(\alpha))$ is a permutation group.

\begin{proposition}
if $(G, \Gamma)$ is a vertex transitive graph
then for all $\alpha, \beta \in V \Gamma$
the permutation groups $(G_\alpha, \Gamma(\alpha))$ and
$(G_\beta, \Gamma(\beta))$ are permutation equivalent.

\begin{proof}
Suppose that $\beta = \alpha^g$. 
Let $\eta : \Gamma(\alpha) \to \Gamma(\beta) $ be given by: $ \gamma \mapsto \gamma^g $,
and let $\varphi : G_\alpha \to G_\beta $
be given by: $ x \mapsto g^{-1} x g $
then:
\begin{eqnarray*}
\eta(\gamma^x) & = & \gamma^{xg} \\
               & = & \gamma^{gg^{-1}xg} \\
	       & = & (\gamma^g)^{\varphi(x)} \\
               & = & \eta(\gamma)^{\varphi(x)}
\end{eqnarray*}	       
for any $\gamma \in \Gamma(\alpha)$ and any $x \in G_\alpha$.
Thus the pair of maps $(\eta,\varphi)$ establish the desired
permutation equivalence.
\end{proof}
\end{proposition}

If $\Gamma$ is $G$-vertex transitive,
then for any permutation group theoretic property $\cP$
we say that $\Gamma$ is {\it $G$-locally $\cP$}
if $(G_\alpha, \Gamma(\alpha))$ has property $\cP$.

\section{Symmetric Graphs}

\begin{definition} [Symmetric Graph]

A symmetric graph is a triple $(G,\Gamma,\rho)$
where $G$ is a group, $\Gamma$ is a graph
and $\rho$ is a homomorphism:
$$\rho : G \to \Aut(\Gamma)$$
such that $\Gamma$ is $G$-vertex transitive and $G$-locally transitive.
\end{definition}

Again, we won't always mention $\rho$ explicitely
and will speak of the symmetric graph $(G,\Gamma)$, or sometimes
the $G$-symmetric graph $\Gamma$.

\begin{definition} [s-arc]
An $s$-arc of a graph $\Gamma$ is a sequence of vertices $(v_0, v_1, ... v_s)$
such that $v_i$ is adjacent to $v_{i+1}$ for each $i$ and $v_{i-1} \neq v_{i+1}$.
The set of $s$-arcs of $\Gamma$ is denoted by $\Arc_s(\Gamma)$.
\end{definition}

If $G$ acts on $\Gamma$ as a group of automorphisms,
then $G$ also acts on $\Arc_s(\Gamma)$ in a natural way:
$$ (v_0, v_1, ... v_s)^g = (v_0^g, v_1^g, ... v_s^g)$$
The graph $\Gamma$ is said to be $(G,s)$-arc transitive if $(G, \Arc_s(\Gamma)$
is transitive.
Historically there has been much interest in highly arc-transitive graphs,
that is graphs that are $s$-arc transitive for large $s$.

One of the oldest results in the area is the theorem by Tutte \cite{TUTTE1,TUTTE2},
proved using combinatorial arguments,
that there are no $s$-arc transitive graphs
of valency $3$ for $s > 5$.
More recently Weiss \cite{WEISS} was able to show, with the aid
of the classification of finite simple groups
that there are no $s$-arc transitive graphs,
of any valency, for $s > 7$.

Symmetric graphs have historically been characterized by their arc transitivity
rather than their local transitivity.
A $(G,0)$-arc transitive graph is just a $G$-vertex transitive graph.
The next theorem shows that $(G,1)$-arc transitive graphs
wihtout isolated vertices are symmetric graphs.

\begin{proposition}
if $(G, \Gamma)$ is a symmetric graph, the $\Gamma$ is $(G,1)$-arc transitive.
Conversly, if $\Gamma$ contains no isolated vertices and is $(G,1)$-arc transitive,
then $(G, \Gamma)$ is a symmetric graph.

\begin{proof}
Suppose that $(G, \Gamma)$ is a symmetric graph, then
$G$ acts on $A \Gamma$ in the obvious way:
$(\alpha,\beta)^x = (\alpha^x, \beta^x)$.
Let $(\alpha_1, \beta_1)$ and $(\alpha_2, \beta_2)$
be any two arcs.
By the $G$-vertex transitivity of $\Gamma$ we can find a $g \in G$ such that
$\alpha_1^g = \alpha_2$.
Since $\beta_1 \in \Gamma(\alpha_1)$ and $g$ is an automorphism,
we must have $\beta_1^g \in \Gamma(\alpha_1^g) = \Gamma(\alpha_2)$.
By $G$-local transitivity we can find $h \in G_{\alpha_2}$
such that $(\beta_1^g)^h = \beta_2$.
It follows that
$(\alpha_1, \beta_2)^{gh} = (\alpha_1^{gh}, \beta_1^{gh})
= (\alpha_2, \beta_2)$.
Thus $G$ acts transitively on the arcs of $\Gamma$.

Now suppose that $\Gamma$ has no isolated vertices and $(G,A \Gamma)$
is transitive. For any two vertices $\alpha_1$ and $\alpha_2$,
there is some arc beginning at $\alpha_1$, say $(\alpha_1,\beta_1)$
and some arc beginning at $\alpha_2$, say $(\alpha_2, \beta_2)$.
By arc transitivity we can find some $g \in G$ carrying
$(\alpha_1, \beta_1)$ to $(\alpha_2, \beta_2)$,
and this $g$ must carry $\alpha_1$ to $\alpha_2$.
So $\Gamma$ is $G$-vertex transitive.

Suppose $\gamma_1$ and $\gamma_2$ are both elements of $\Gamma(\alpha)$.
The $(\alpha,\gamma_1), (\alpha,\gamma_2) \in A \Gamma$.
By arc transitivity we can find a $g \in G$ such that
$(\alpha, \gamma_1)^g = (\alpha, \gamma_2)$
and this $g$ must carry $\gamma_1$ to $\gamma_2$.
Thus $\Gamma$ is $G$-locally transitive.
The result follows.
\end{proof}
\end{proposition}

\section{Orbital Graphs}

If $(G,\Omega)$ is a transitive permutation group,
then $G$ has a natural action on $\Omega \times \Omega$
given by:
$$ (\alpha,\beta)^g = (\alpha^g, \beta^g) $$
Although $(G,\Omega)$ is transitive, $(G,\Omega \times \Omega)$
need not be. 
The orbits of $G$ on $\Omega \times \Omega$ are called {\it orbitals}

\begin{lemma} 
If $(G,\Omega)$ is a transitive permutation group
and $\alpha \in \Omega$
then there is a natural correspondence between the orbitals of $(G,\Omega)$
and the orbits of $(G_\alpha,\Omega)$ 
\begin{proof}
Let $\Delta$ be any orbital of $(G,\Omega)$
and let $(\gamma,\delta)$ be any element of $\Delta$.
Let $g \in G$ be such that $\gamma^g = \alpha$.
Since $\Delta$ is an orbital, it follows that
$(\alpha,\delta^g) \in \Delta$.
Thus every orbital of $(G,\Omega)$ contains an element
of the form $(\alpha,\beta)$.

Let $\eta : \Omega \to \Omega \times \Omega$ be the map given by
$\beta \mapsto (\alpha, \beta)$. 
Suppose that $\beta_1$ and $\beta_2$
lie in the same orbit of $(G_\alpha,\Omega)$.
Then there is some $g \in G_\alpha$ such that $\beta_1^g = \beta_2$.
It follows that $(\alpha,\beta_1)^g = (\alpha,\beta_2)$
and so $(\alpha,\beta_1)$ and $(\alpha,\beta_2)$ lie in the same
orbital of $(G,\Omega)$.

Conversely suppose that $(\alpha,\beta_1)$ and $(\alpha,\beta_2)$
lie in the same orbital of $(G,\Omega)$.
Then there is some $g \in G$ such that $(\alpha,\beta_1)^g = (\alpha,\beta_2)$.
That is $\alpha^g = \alpha$ and $\beta_1^g = \beta_2$.
It follows that $g \in G_\alpha$ and $\beta_1$ and $\beta_2$
lie in the same orbit of $(G_\alpha, \Omega)$.
Thus, the map $\eta$ in fact induces a map from the orbits of $(G_\alpha,\Omega)$
to the orbitals of $(G,\Omega)$.
\end{proof}
\end{lemma}

The orbital $\{ (\alpha,\alpha) : \alpha \in \Omega)$
is given a special name. Its called the {\it diagonal} orbital.
The {\it rank} of a permutation group is the number of orbitals.
An orbital $\Delta$ is said to be {\it self-paired} if
$$ (\alpha, \beta) \in \Delta \Longleftrightarrow
   (\beta, \alpha) \in \Delta $$

\begin{definition} [Orbital Graph]
If $(G,\Omega)$ is a transitive permutation group and
$\Delta$ is a self-paired orbital 
then the {\it orbital graph} $\orb_\Delta(G, \Omega)$
is the graph with vertex set $\Omega$ and arc-set $\Delta$
Need to say more about the diagonal orbital.
\end{definition}

Note that if $\Delta$ is taken to be the diagonal orbit,
then the resulting graph has loops, and so is not actually
a graph by our definition. Generally we assume that $\Delta$
is not the diagonal orbit.

\begin{proposition}
Orbital graphs are symmetric. Every symmetric graph is an orbital graph
\begin{proof}
Suppose that $\Gamma = \orb_\Delta(G,\Omega)$ 
for some transitive permutation group
$(G,\Omega)$ and some self-paired orbital $\Delta$.
Since $V \Gamma = \Omega$ and $(G,\Omega)$ is transitive,
$\Gamma$ is $G$-vertex transitive.
Since $A \Gamma = \Delta$ and $\Delta$ is an orbital
$\Gamma$ is $G$-arc transitive.
It follows that $\Gamma$ is symmetric.

Now suppose that $\Gamma$ is a $G$-symmetric graph.
Since $\Gamma$ is $G$-vertex transitive
$(G, V \Gamma)$ is a transitive permutation group.
Since $A \Gamma \subseteq V \Gamma \times V \Gamma$
and $\Gamma$ is $G$-arc transitive, $A \Gamma$
is an orbital of $(G,V \Gamma)$.
It follows that $\Gamma = \orb_{A \Gamma}(G,V \Gamma)$.
\end{proof}
\end{proposition}

Relaxing the condition that $\Delta$ must be self-paired 
leads to a symmetric digraph.

\section{Imprimitive Symmetric Graphs}

Recall that if $\Delta$ is a subset of $\Omega$ 
and $g$ is an element of $G$, then $\Delta^g$ denotes
the image of $\Delta$ under the action of $g$.

\begin{definition}[Block of Imprimitivity]
If $(G,\Omega)$ is a transitive permutation group
then a subset $\Delta$ of $\Omega$
is said to be a \emph{block of imprimitivity} if for every $x \in G$
either $\Delta^x = \Delta$ or $\Delta^x \cap \Delta = \emptyset$
\end{definition}

\begin{definition}[$G$-invariant partition]
If $(G, \Omega)$ is a transitive permutation group,
then a partition $\cB$ of $\Omega$ is said to be \emph{$G$-invariant}
if for each $\Delta \in {\cB}$ and each $x \in G$
we have $\Delta^x \in {\cB}$. 
That is, $\cB$ admits $G$ as a group of permutations in a natural way.
\end{definition}

\begin{proposition}
If $(G,\Omega)$ is a transitive permutation group
and $\Delta$ is a block of imprimitivity
then $\cB = \{ \Delta^g : g \in G \}$ is a $G$-invariant
partition.
\begin{proof}
Since $(G,\Omega)$ is transitive,
every $\alpha \in \Omega$ is contained in $\Delta^g$
for some $g \in G$. Thus $\bigcup_{g \in G} \Delta^g = \Omega$.
If $\Delta^{g_1} \cap \Delta^{g_2} \neq \emptyset$
for some $g_1, g_2 \in G$
then $\Delta \cap \Delta^{g_2 g_1^{-1}} \neq \emptyset$.
Since $\Delta$ is a block of imprimitivity,
this implies that $\Delta = \Delta^{g_2 g_1^{-1}}$
and so $\Delta^{g_1} = \Delta^{g_2}$.
Thus $\cB$ is a partition of $\Omega$.
For any $\Delta^g \in \cB$ and any $x \in G$
$(\Delta^g)^x = \Delta^{gx} \in \cB$.
Thus $\cB$ is a $G$-invariant partition of $\Omega$.
\end{proof}
\end{proposition}

\begin{proposition}
If $(G,\Omega)$ is a transitive permutation group
and $\cB$ is a $G$-invariant partition of $\Omega$,
then each $B$ in $\cB$ is a block of imprimitivity.
\begin{proof}
Since $\cB$ is $G$-invariant, $B^g \in \cB$ for any
$g \in G$. Since $\cB$ is a partition, if $B^g \neq B$
then $B^g \cap B = \emptyset$.
The result follows.
\end{proof}
\end{proposition}

\begin{proposition}
The map $\pi : \Omega \to \cB$ which sends each point of $\Omega$
to the block of $\cB$ containing it
induces a surjective permutation homomorphism from $(G, \Omega)$
to $(G,\cB)$.
\begin{proof}
Need to show that $\pi(\alpha^x) = \pi(\alpha)^x$
for any $\alpha \in \Omega$ and any $x \in G$.
By the previous proposition $\pi(\alpha)$ is a block of imprimitivity.
Let $\Delta = \pi(\alpha)$. 
Since $\alpha \in \Delta$ we must have $\alpha^x \in \Delta^x$
That is $\pi(\alpha^x) = \Delta^x = \pi(\alpha)^x$.
\end{proof}
\end{proposition}

For any $B \in \cB$ we refer to the set $\pi^{-1}(B)$ as the
{\it fiber} of the homomorphism $\pi$ at $B$.
The fibers of a permutation homomorphism are blocks of imprimitivity.
Conversly each block of imprimitivity $\Delta$ 
gives rise to the $G$-invariant partition:
$\cB = \{ \Delta^x : x \in G \} $
and thus also to the homomorphism $(G,\Omega) \mapsto (G,\cB)$.


For every permutation group both $\{a\}$ and $\Omega$ are 
trivially blocks of imprimitivity. A permutation group is \emph{primitive}
if it admits no nontrivial blocks of imprimitivity.
%In category theoretic terms, the primitive representations
%of a group $G$ are those objects of $G$-Set having no proper quotients objects.
%Thus we see that the notions of transitivity and primitivity
%are category theoretic duals.

\begin{definition} [Imprimitive Symmetric Graph]
A symmetric graph $(G,\Gamma)$ is said to be {\it imprimitive} if
the induced permutation group, $(G, V \Gamma)$, is imprimitive.
\end{definition}

\begin{definition} [Quotient Graph]
Suppose that $(G,\Gamma)$ is an imprimitive symmetric graph
with $\cB$ a nontrivial $G$-invariant partition of the vertex set $V \Gamma$.
Define the {\it quotient graph} $\Gamma_\cB$ of $(G,\Gamma)$ with respect to $\cB$
to be the graph with vertex set $\cB$, and an arc $(B,C)$,
whenever there is some $\alpha \in B$ 
and some $\beta \in C$ 
such that  $(\alpha, \beta )$ is an arc of $\Gamma$.
\end{definition}

\begin{proposition}
$(G, \Gamma_\cB)$ is a symmetric graph.

\begin{proof}
We show that $G$ acts transitively on the arcs of $\Gamma_\cB$.
Let $(B,C)$ and $(D,E)$ be any two arcs of $\Gamma_\cB$.
By the definition of the quotient graph,
there must be some $\alpha \in \pi^{-1}(B)$ 
and $\beta \in \pi^{-1}(C)$
such that $(\alpha, \beta) \in A \Gamma$.
Simmilarly, there must be some $\gamma \in \pi^{-1}(D)$ 
and $\delta \in \pi^{-1}(E)$
such that $(\gamma, \delta) \in A \Gamma$,
By the symmetry of $(G, \Gamma)$, we can find an element $g \in G$
which carries  $(\alpha,\beta)$ to $(\gamma,\delta)$.
Since $\pi$ is a $G$-homomorphism, it follows that:
\begin{eqnarray*}
(B,C)^g & = & (\pi(\alpha),\pi(\beta))^g \\
        & = & (\pi(\alpha)^g,\pi(\beta)^g) \\
	& = & (\pi(\alpha^g),\pi(\beta^g)) \\
        & = & (\pi(\gamma),\pi(\delta)) \\
	& = & (D,E).
\end{eqnarray*}
\end{proof}
\end{proposition}

Suppose that some arc of $\Gamma$ 
has both its endpoints in the same fiber of $\cB$,
that is, there exists some $(\alpha,\beta) \in A \Gamma$
such that $\pi(\alpha) = \pi(\beta)$.
Since the fibers are blocks of imprimitivity,
it follows that $\pi(\alpha^g) = \pi(\beta^g)$ for each $g \in G$,
The arc transitivity of $\Gamma_\cB$ then implies 
every arc of $\Gamma$ has both its endpoints in the same fiber. 
In this case $\Gamma_\cB$ is the empty graph (no edges) 
with one vertex per connected component of $\Gamma$. 
This case is not very interesting,
to exclude it we say that the quotient is {\it nontrivial}
if it has valency at least one.
By the above discussion, $\Gamma_\cB$ is nontrivial
if and only if each of the fibers of $\cB$ is an
{\it independent set} of $\Gamma$.
For the remainder of this paper we shall always assume
that the quotient of a $G$-symmetric graph homomorphism
is non-trivial, even if we neglect to state this explicitely.

%Butterflies in bloom. See how the grow from a small seed into such majestic beings?

